\subsection*{Assembler Source Grammar}
The assembler source language is a sequence of newline-terminated statements. Whitespace may separate tokens except
where an instruction spelling joins a mnemonic, size or condition suffix, or modifier. A label may precede a
statement on the same line, and a line may contain labels without an instruction or directive.

\begin{verbatim}
source              ::= { line }
line                ::= { label ":" } [ statement ] end-of-line
statement           ::= instruction | directive
instruction         ::= instruction-form [ ":" "LEN" integer
                        [ "," padding-byte-list ] ]
instruction-form    ::= mnemonic-form [ operand-list ]
operand-list        ::= operand { "," operand }
padding-byte-list   ::= padding-byte { "," padding-byte }
padding-byte        ::= integer | "ILLEGAL" | "NOP"
operand             ::= register-name
                      | expression
                      | address-expression
                      | register-mask
                      | "(" instruction ")"
address-expression  ::= "[" ea-expression "]"
symbol-reference    ::= bare-identifier | quoted-identifier
quoted-identifier   ::= "`" { quoted-character | escape } "`"
escape              ::= "\`" | "\\"
\end{verbatim}

This grammar is an envelope, not a second declaration of the instruction set. The instruction encoding catalog is
the sole definition of the accepted mnemonic forms, suffixes, modifiers, operand counts, operand ordering, and
operand classes. A parenthesized instruction is accepted only where an instruction form declares an instruction
operand. The effective-address mode catalog is the sole definition of the accepted \texttt{ea-expression} forms.
The brackets are part of every address expression and never form a non-address operand. Whether a consuming
instruction reads memory, writes memory, computes the address only, or applies another address-directed operation is
defined by that instruction\textquotesingle{}s operand role.

The optional \texttt{: LEN n} annotation requests an encoded extended instruction length of \texttt{n} bytes. If no
padding-byte list follows, the assembler fills every trailing (:term:base.terms.padding:) byte with (:ref:base.instructions.ILLEGAL:), whose canonical
one-byte encoding is zero. If a list is present, each item specifies one trailing (:term:base.terms.padding:) byte in instruction-stream
order; its item count must equal the requested length minus the (:term:base.terms.required_instruction_length:). (:ref:base.instructions.NOP:) denotes its
canonical one-byte encoding, and an integer (:term:base.terms.padding:) byte must be an absolute value from 0 through 255. Padding bytes do
not accept expressions or relocations.

Register names are recognized case-insensitively. Symbol names are case-sensitive. A bare identifier that is a
register name denotes that register in a context that accepts it; every other bare identifier in an expression
denotes a symbol. Backtick quoting may be applied to any symbol reference and forces the enclosed name to
be interpreted as a symbol, so \texttt{R1} denotes the register while \texttt{`R1`} denotes the distinct symbol named
\texttt{R1}. Within a quoted identifier, \texttt{\textbackslash\textasciigrave} denotes a literal backtick and
\texttt{\textbackslash\textbackslash} denotes a literal backslash. No other escape is defined, and a quoted
identifier cannot cross a newline.

A relocation annotation follows the complete symbol reference and is not part of the quoted name. For example,
\texttt{`R1`@abs64} applies \texttt{abs64} to the symbol \texttt{R1}, whereas \texttt{`R1@abs64`} names a symbol
whose name contains the characters \texttt{@abs64}. An assembler emits an unquoted symbol when its name is a valid
bare identifier that does not collide with a register name; otherwise it emits the canonical backtick-quoted
form.
